% Chapter 10 -- Applications (Part II, Reference).
% Sources (Law 1): os/flair/tbxgate.c (the Toolbox Gate: one disk
% application resident at a time, TENANT-SLOT-BUSY; a failing application
% is closed, TENANT-CRASH, and its windows removed); os/flair/process.c
% (the foreground application; its bar in band 2); os/milton/kmain.c
% (launch by double-click of an .EXE, flair_live_launch_app; the Apple-menu
% launch; the close box of a window made to ask; tbx_ask_quit before the
% power commands; sf_run: the Open and Save dialog boxes, sf_name_ok, the
% Replace question); os/flair/scrap.c / scrap.h (one desk scrap in memory,
% TEXT up to 8192 bytes, kept until the computer is restarted);
% os/flair/stdfile.c; os/apps/itext/itext.c (the one user of the Scrap in
% this build).  Model: Mac User's Guide 1991 Ch. 11 (Using programs) and
% the System 6 guide Ch. 10 MultiFinder (apple-macintosh.md Sec 8.1, 8.2).
% GATES: test-flair-app-launch test-tenant-load test-flair-tenant-menu test-tbx-menu test-scrap test-stdfile test-flair-itext test-flair-files test-flair-files-reap test-pubs-flair-screens
\Gates{test-tenant-load, test-flair-tenant-menu, test-flair-i123,
  test-flair-iword,
  test-tbx-menu, test-scrap, test-stdfile, test-flair-itext,
  test-flair-files, test-flair-files-reap, test-pubs-flair-screens}
\chapter{Applications}
\label{ch:appsref}

\section{Two kinds of application}
\index{applications}
\SideHead{Built in}The Finder, HELLO and NOTES are part of FLAIR and are
running when it starts.  The Finder cannot be quit; HELLO and NOTES end when
you close their windows.

\SideHead{From the disk}All other applications are programs on the data
disk, in the \file{APPS} folder, with names ending in \file{.EXE}.  FLAIR
loads one when you open it, and frees the memory it used when it quits.
Appendix~A lists the applications on the data disk.

\section{Opening and quitting}
\index{applications!opening}\index{applications!quitting}
Open an application from the disk by double-clicking its icon, or, for an
accessory, by choosing it from the Apple menu.  One application from the
disk can be open at a time; while one is open, opening another does
nothing.

An application quits when you choose \menuitem{Quit} (or, in some
applications, \menuitem{Close}) from its \menuname{File} menu, or click the
close box of its window.  An application that holds work you have not
saved asks you first whether to save it.

When you choose \menuitem{Restart}, \menuitem{Shut Down} or \menuitem{Exit
to InitechDOS}, FLAIR asks the open application to quit, exactly as if you
had chosen its \menuitem{Quit} command.  If the application asks about
unsaved work and you click \texttt{Cancel}, the command stops.

If an application stops working, FLAIR closes it: its windows disappear,
its menus leave the menu bar, and the rest of the desktop goes on working.
Work it had not saved is lost.

\section{Menus and the keyboard}
\index{menu bars!lower}
While an application is active, the lower menu bar shows its menus, and
the \key{Ctrl} key combinations shown in them choose its commands.  An
application's menus are its own: one may have only a \menuname{File} menu,
another a full set.  Most applications from the disk also have the Apple
menu at the left of their menu bar.

\section{The Clipboard}
\index{Clipboard}\index{Scrap}
When you choose \menuitem{Cut} or \menuitem{Copy} in an application, the
selected text is placed on the \term{Clipboard}, a holding place in the
computer's memory, replacing what was there.  \menuitem{Paste} puts a copy
of what is on the Clipboard where you are working.  The Clipboard holds
text, up to 8,192 characters, and keeps it when the application quits, so
you can paste it into a document you open later.  It is emptied when the
computer is restarted; it is never written to the disk.

In this version of FLAIR, InitechText uses the Clipboard.  The
\menuname{Edit} menus of the Finder and the Calculator are dimmed.

\section{The Open and Save dialog boxes}
\index{dialog boxes!Open}\index{dialog boxes!Save}
Applications that work with documents use the same two dialog boxes to ask
which document to open and where to save one.  Figure~\ref{fig:textsaveas}
and Figure~\ref{fig:textopenbox} in Chapter~\ref{ch:applications} show them.

\SideHead{Folder menu}The box at the top names the folder whose contents
are listed.  Click it to go up to the folder that holds it.

\SideHead{List}The files and folders in that folder, one to a line.  Click
a name to select it; move the selection with \key{Up Arrow} and \key{Down
Arrow}.  Double-click a folder to list its contents; in the Open box,
double-click a document to open it.

\SideHead{Disk name}Above the buttons, the name of the disk:
\texttt{INITECH}.

\SideHead{Buttons}\texttt{Open} or \texttt{Save}, the default button, does
what the box is for.  \texttt{Cancel} puts the box away and does nothing.
\texttt{Desktop} lists the top level of the disk.  \texttt{Eject}, and in
the Save box \texttt{New}, are dimmed.

\SideHead{Name box}In the Save box, the name the document will be saved
under.  Type the name; letters appear in capitals.  Press \key{Backspace}
to erase.  Typing a name deselects the list.

A name follows the InitechDOS rules: one to eight characters, then, if
you want, a period and one to three more.  The Save box does not accept
any other name.  If the name is already used in the folder, an alert asks
whether to replace that file; if it is the name of a folder, nothing is
saved.

\section{Initech 123 and InitechWord}
\index{Initech 123}\index{InitechWord}
The data disk also holds two applications that are sold and described
separately: Initech 123, the spreadsheet (\file{123.EXE}), and InitechWord,
the word processor (\file{WORD.EXE}).  Each opens from its icon in the
\file{APPS} folder like any other application from the disk, in a window
that fills the screen below the menu bars.  Its \menuname{File} menu is in
the lower menu bar, and the rest of its commands are given from the
keyboard, as its own books describe.

\FlairScreen[label=fig:o123]{o_123}{Initech 123, opened from the APPS
folder}

\FlairScreen[label=fig:oword]{o_word}{InitechWord, opened from the APPS
folder}

While one of them is open, it is the one application from the disk: quit it
before you open another (Chapter~\ref{ch:multi}).  Initech 123 needs a
math coprocessor.  \doctitle{Initech 123 Reference} and
\doctitle{InitechWord Reference} describe the two applications.

\section{Programs that run at the InitechDOS prompt}
\index{InitechDOS programs}
The \file{APPS} folder also holds programs whose names end in \file{.COM}:
the InitechDOS utility programs (\file{ATTRIB}, \file{CHKDSK},
\file{EDLIN}, \file{LABEL}, \file{TREE}, \file{XCOPY}) and the TPS Report
Generator (\file{TPSREP}).  They are not opened from the desktop.  Leave
FLAIR with \menuitem{Exit to InitechDOS}; at the prompt, type \type{CD APPS}
to change to the folder, then type the program's name.  The \doctitle{InitechDOS Reference} describes the utility
programs; the TPS Report Generator has its own book.
